% !TeX root = ../papers/nc-p3.tex
\section{周期的に続く安定な文字作用の一意性}
\label{sec:periodic-stable-arrow-uniqueness}
\begin{theorem}\label{thm:periodic-stable-arrow-uniqueness}
二次四窓が多項式成長するとする。同じ次元 $d$ の空間の周期列で、各辺の四矢 $A_i$ が
全体を次の全体へ送り、各非零真部分空間 $P$ で $\dim A_i(V_i)P>\dim P$ を満たすものを考える。
周期は4の任意の倍数でよく、辺ごとにこれを二頂点四矢のKronecker表現として読む。
二つの両立する周期列で、同じ位置の一辺のKronecker表現が同型なら、
両者は共通周期に展開したとき頂点基底変更と共通次数倍率を除いて同じ列である。
特に有限四辺では、全位置が二次元の周期列はどの一つの向き付き行列式二次曲面からも
周期列の中で一意に決まる。辺単独の許す後続が網である場合も含む。
\end{theorem}
\begin{proof}
共通周期 $n$ に取り直す。$n$ 頂点折り畳みの各始点の各次数の経路空間は、四頂点の場合と同じなので
多項式成長が保存される。与えた二つの表現 $S,T$ は単純である。部分表現の次元列は非減少で、
非零真部分空間があれば真に増えるからである。一辺を消した切断は
重み $(1,\ldots,1,-(n-1))$ に関して安定である。非減少列の最後が前の全て以上で、
全て等しい真部分次元列は厳密な増加条件に反する。従って切断の自己準同型はスカラーで、
非同型な二つのこうした切断の間のHomは零である。

共通辺の次の辺だけを可変にして切断する。切断が同型なら定理\ref{thm:simple-linear-completion-growth}の
一辺版が共通次数倍率までの同型を与える。その証明は $n$ 頂点でも同じである。
切断が異なると仮定し、共通辺の両端の基底を選んでその行列を一致させる。
$S\to T$ の非対角な一次変形として、可変辺で $T$ 自身の矢を使い、残りを零とする。
前の関係は共通辺が同じなので満たし、次の関係は $T$ の関係そのものである。
これが内部変形なら可変辺を除いた切断間の準同型から来るはずで、そこは零である。
逆方向にも $S$ の矢により非自明な変形がある。さらに $S\to S$ に同じ辺だけの倍率方向がある。
非零な正長閉路のtraceの微分によりこれは内部でない。

この三方向をそれぞれ二頂点 $a,b$ の経路環の矢 $y:a\to b,\ z:b\to a,\ x:a\to a$ で係数づけ、
定数部を $S\oplus T$ としたブロック行列表現を作る。対角の基礎体は対応する頂点冪等元で置く。
全ての可変項が一辺に集中するため二次関係は非可換経路係数でも厳密に成立する。
矢イデアルの自乗を割った像を $E$ とすると、その定数像は
$D=\operatorname{End}_k(S)\oplus\operatorname{End}_k(T)$ の全体である。
非同型単純の核は相異なる極大両側イデアルであり、Burnsideと中国剰余定理による。
定数への核は $D$ 両側加群で、上の三方向の矩形行列成分はいずれも非同型単純な $D$ 両側加群である。
三成分が全て核に含まれる。一成分を欠けば、他の二成分を捨てた像は $D$ のコピーで、
その一次項は $D$ から矩形行列への導分となる。有限分離代数の導分は内部なので、対応する一次変形の
非自明性に反する。各成分の $(1,1)$ 行列単位にその矢を掛けたものを元の経路代数に持ち上げる。
$a$ におけるループ $x$ と往復 $yz$（合成順は経路に従う）の任意の語は、
最低次で対応する経路と非零行列単位を与える。この二つの語の連接は一意に読めるので
自由二生成代数が得られ、多項式成長に反する。
二次元では厳密な増加は各非零ベクトルでの全射であり、一辺のKronecker同型を
行列式と滑らかな場合の罫線で判定できることは定理\ref{thm:finite-two-dimensional-quadrics}による。
\end{proof}

\begin{corollary}\label{cor:quadrics-period-fold-simple}
有限四辺の多項式成長窓において、向き付き二次曲面の周期列の最小周期を $4m$
（位置を保存して測る）とする。この列の単純を四頂点に押し下げると各頂点次元 $2m$ の単純で、
一周の次数倍率の安定化群はちょうど $\mu_m$ である。
各位置の $m$ 個の向き付き二次曲面は全て異なる。一隅が $s$ 行 $s$ 列行列の恒等式を満たす場合は
$m\le\lfloor s/2\rfloor$ である。
\end{corollary}
\begin{proof}
同じ位置で同じ向き付き曲面が二度出れば元の列とその移動に定理を使え、より短い周期を得る。
逆に射影データの短い周期は接続線の一意性により頂点変更と倍率を除いた列の周期である。
四頂点へ押し下げた表現 $U$ を $4m$ 頂点へ引き戻すと、全ての四段移動の直和に分かれる。
各位置で $m$ 個の元の頂点空間を並べ、戻り矢が巡回することを見ればよい。
これらは各位置の曲面が違うので非同型単純で、部分表現はその幾つかの直和に限る。
$U$ の部分表現の引戻しは移動に不変で、全ての直和因子が推移的に移るため零か全体である。
$U$ を一辺の倍率 $\lambda$ で捩ったものの引戻しでは、各直和因子の長い一周の倍率が $\lambda^m$ 倍される。
曲面が違うので同型なら同じ因子へ戻る必要があり、そこでは長い切断の自己射がスカラーなので
$\lambda^m=1$ が必要である。逆にこの条件なら $m$ 個の層の倍率で捩りを吸収できる。
隅の単純作用が $M_{2m}$ 全体なので行列恒等式の上界も従う。
\end{proof}
